% Generated draft: requires mathematical and source review.
\documentclass[11pt]{article}
\usepackage[margin=25mm]{geometry}
\usepackage{fontspec}
\setmainfont{Latin Modern Roman}
\usepackage{amsmath,amssymb,mathtools,unicode-math}
\setmathfont{Latin Modern Math}
\usepackage{enumitem,needspace,xurl,hyperref}
\setlength{\parindent}{0pt}
\setlength{\parskip}{5pt}
\setlength{\emergencystretch}{3em}
\newcommand{\N}{\mathbb{N}}
\newcommand{\Z}{\mathbb{Z}}
\newcommand{\Q}{\mathbb{Q}}
\newcommand{\R}{\mathbb{R}}
\newcommand{\C}{\mathbb{C}}
\newcommand{\F}{\mathbb{F}}
\newcommand{\A}{\mathbb{A}}
\newcommand{\PP}{\mathbb{P}}
\newcommand{\E}{\mathbb{E}}
\newcommand{\eps}{\varepsilon}
\newcommand{\dd}{\,\mathrm{d}}
\begin{document}
\section*{UnsolvedMath ID: 2305049}
\textbf{Source problem code:} catalogue entry
\subsection{Definitions and mathematical statement}
In the weight-zero function convention, let a group $\Gamma$ act by conformal transformations on a plane domain $D$, and let $f$ be a nonconstant meromorphic function on $D$ with $f(\gamma z)=f(z)$ for all $\gamma\in\Gamma$ wherever defined. At a point $z_0$ where $f$ is holomorphic, write its Taylor expansion in a specified coordinate as
\[
f(z)=\sum_{n=0}^\infty a_n(z-z_0)^n.
\]
Its gap pattern is the set of indices at which $a_n=0$, or equivalently its nonzero coefficient support. If that support is infinite, list its positive indices as $n_1<n_2<\cdots$. A Hadamard-gap expansion means that for some $q>1$,
\[
n_{j+1}/n_j\geq q
\quad\text{for all sufficiently large }j.
\]
Finite Taylor polynomials are a separate case, not an infinite Hadamard-gap sequence.

L. A. Rubel asks which gap patterns are possible for nonconstant automorphic functions, in particular whether a Hadamard-gap expansion can occur. The source expects a sharp answer to depend on the automorphy group. It does not specify the group class, domain, expansion center, coordinate, or whether a weight or character is allowed instead of weight-zero invariance. Thus a universal exclusion or an existence claim for a particular group cannot be inferred from the question alone; the definitions above make one standard interpretation explicit.

The source relates the question to a theorem attributed to R\'enyi, with ``no citation,'' saying that a nonconstant periodic entire function cannot have more than half of its Taylor coefficients zero. The meaning of ``more than half'' is not defined (for example, which asymptotic density or finite counting statement is intended), so no precise coefficient-density inequality is asserted here. The supplied summary reports the automorphic-gap question open as of the 2018 edition, with no recent resolution located.
\subsection{Short English statement}
Rubel asks how sparse the Taylor coefficients of an automorphic function can be, especially whether their nonzero indices can grow geometrically. The answer is expected to depend on the group, which is unspecified. A related periodic-function coefficient claim is quoted without a citation or precise density convention.
\subsection{Sources}
\begin{itemize}
\item Local UnsolvedMath registry, permanent ID 2305049, exact source code catalogue entry, ``Research Problems in Function Theory --- Problem 5.49''; status open, classification OPEN-TRIAGE.
\item Hayman, \emph{Research Problems in Function Theory} (2018), Problem 5.49, attributed to L. A. Rubel, \url{https://arxiv.org/abs/1809.07200}. The record identifies extraction from public arXiv source TeX, accessed July 29, 2026; the 2018 update was retained for status triage.
\item R\'enyi, periodic-entire-function coefficient theorem, explicitly marked ``no citation'' in the supplied source. No precise density statement or authoritative bibliographic reference is supplied.
\item The supplied research summary reports openness as of the 2018 edition and no recent resolution located, with no further citation or precise search date.
\end{itemize}
\end{document}
